\documentclass[11pt]{article}
\usepackage[a4paper,margin=28mm]{geometry}
\usepackage{amsmath,amssymb,amsthm,mathtools}
\usepackage[colorlinks=true,linkcolor=blue,urlcolor=blue]{hyperref}
\newtheorem{definition}{Definition}[section]
\newcommand{\C}{\mathbb C}
\newcommand{\ii}{\mathrm i}
\newcommand{\ind}{\operatorname{ind}}
\newcommand{\coeff}{\operatorname{coeff}}
\title{Local multiplicities at balanced singular configurations\\
\large Definitions and the alternating local index}
\author{}
\date{November 1, 2026}
\begin{document}
\maketitle
\vspace{-2em}

This note sets up the local multiplicity question for the auxiliary
cycles $Z_{A,B;\pi}$ of \texttt{../main.tex}.
We use its rescaled coordinates, in which the singular values are
$1,-1$ (equivalently, $\lambda=\ii x/2$ has singular values
$\ii/2,-\ii/2$). We stop at the definition of the alternating index.
In particular, no interpretation of this index as a count of deformed
components is assumed.

\section{The polynomial equations and the point}

Fix integers $N\geq 2$ and $L\geq 2N$, and write $[N]=\{1,\ldots,N\}$.
The ordered Bethe polynomials are
\begin{equation}\label{eq:bethe}
 P_i(x)=(x_i+1)^L\prod_{j\ne i}(x_i-x_j-2)
       -(x_i-1)^L\prod_{j\ne i}(x_i-x_j+2),
 \qquad i\in[N].
\end{equation}
Fix $1\leq k\leq\lfloor N/2\rfloor$ and put $r=N-2k$. Our point is
\[
 p=(\underbrace{1,\ldots,1}_{k},
    \underbrace{-1,\ldots,-1}_{k},z_1,\ldots,z_r),
 \qquad z_u\ne\pm1,\quad z_u\ne z_v\ (u\ne v).
\]
The notation $1^k,-1^k$ means repeated coordinates, not powers.
Its singular label sets are
\[
 A_0=\{1,\ldots,k\},\qquad B_0=\{k+1,\ldots,2k\},
 \qquad R_0=[N]\setminus(A_0\cup B_0).
\]
All coordinates and multiplicities below are ordered; we do not pass
to permutation orbits.

\section{Fixing singular coordinates and imposing a partition}

Take disjoint subsets $A,B\subseteq[N]$, where either both are empty
or both are nonempty. Set $a=|A|$, $b=|B|$, and
$R=[N]\setminus(A\cup B)$. For $t\in R$, define
\begin{equation}\label{eq:reduced}
\begin{aligned}
 P_t^{A,B}(x_R)
 &=
 \frac{P_t(x)|_{x_\alpha=1\ (\alpha\in A),\,x_\beta=-1\ (\beta\in B)}}
      {(x_t+1)^a(x_t-1)^b}\\
 &=(x_t+1)^{L-a}(x_t-3)^a
       \prod_{\substack{u\in R\\u\ne t}}(x_t-x_u-2)\\
 &\quad -(x_t-1)^{L-b}(x_t+3)^b
       \prod_{\substack{u\in R\\u\ne t}}(x_t-x_u+2).
\end{aligned}
\end{equation}
The quotient is a polynomial identity, including at $x_t=\pm1$.
Thus this definition removes a common polynomial factor;
it is not division of numerical values at $p$.

Let $\pi=\{S_1,\ldots,S_m\}$ be a set partition of $R$, and choose
$t_\alpha\in S_\alpha$. Define
\begin{equation}\label{eq:ideal}
\begin{aligned}
 I_{A,B;\pi}=\bigl(&x_\alpha-1\ (\alpha\in A),\
                    x_\beta+1\ (\beta\in B),\\
                 &P_{t_\alpha}^{A,B}\ (1\leq\alpha\leq m),\\
                 &x_u-x_{t_\alpha}\
                   (u\in S_\alpha\setminus\{t_\alpha\},
                    \ 1\leq\alpha\leq m)\bigr).
\end{aligned}
\end{equation}
After imposing the block equalities, equations indexed by labels in
one block agree. The ideal is therefore independent of the chosen
representatives. There are $N$ displayed generators in total.
Write $\hat0_R$ for the partition of $R$ into singletons.

Equivalently, replace all $x_u$ in $S_\alpha$ by a single variable
$y_\alpha$ in $P_{t_\alpha}^{A,B}$. If the resulting polynomial is
$H_\alpha(y)$, the quotient algebra in \eqref{eq:ideal} is
\[
 \C[x_1,\ldots,x_N]/I_{A,B;\pi}
 \simeq \C[y_1,\ldots,y_m]/(H_1,\ldots,H_m).
\]
The fixed coordinates and block equalities have already been eliminated
in this presentation.

\section{Local multiplicity and the meaning of a simple zero}

Assume the auxiliary ideals used below are zero-dimensional.
For a point $q$, let $\mathfrak m_q$ be its maximal ideal and define
\begin{equation}\label{eq:local-length}
 m_{A,B;\pi}(q)=
 \dim_\C\left(\C[x_1,\ldots,x_N]/I_{A,B;\pi}\right)_{\mathfrak m_q},
\end{equation}
with value zero if $q\notin V(I_{A,B;\pi})$.
The associated cycle is
\[
 Z_{A,B;\pi}=\sum_q m_{A,B;\pi}(q)[q].
\]
Thus the multiplicity of $p$ in $Z_{A,B;\pi}$ means exactly the
integer $m_{A,B;\pi}(p)$.

For the point above, fix all its singular coordinates and use the
singleton partition on the ordinary labels. The remaining equations,
in variables $y_1,\ldots,y_r$, are
\begin{equation}\label{eq:fully-reduced}
\begin{aligned}
 F_u(y)&=(y_u+1)^{L-k}(y_u-3)^k
                  \prod_{v\ne u}(y_u-y_v-2)\\
       &\quad -(y_u-1)^{L-k}(y_u+3)^k
                  \prod_{v\ne u}(y_u-y_v+2),
 \qquad 1\leq u\leq r.
\end{aligned}
\end{equation}
We assume $F_u(z)=0$ for every $u$. This implies that $p$ solves
the original equations \eqref{eq:bethe}: the singular equations vanish
automatically, and the removed factors at ordinary coordinates are
nonzero.

\begin{definition}[Simple fully reduced zero]\label{def:simple}
The configuration $p$ is called simple for the fully reduced system if
\[
 \det\left(\frac{\partial F_u}{\partial y_v}(z)\right)_{u,v=1}^r
 \ne0.
\]
Equivalently,
$m_{A_0,B_0;\hat0_{R_0}}(p)=1$.
For $r=0$, the system is empty, the determinant is $1$ by convention,
and its local algebra is $\C$.
\end{definition}

This is the precise meaning of a simple zero used in this note.
It concerns the ordinary coordinates after \emph{all} $k$ pairs have
been fixed. It does not assert that $p$ is a simple zero of the full
Bethe ideal or of every $I_{A,B;\pi}$.

For a particular compatible partition ideal, its multiplicity can be
computed instead as
\[
 m_{A,B;\pi}(p)
 =
 \dim_\C
 \frac{\C[[u_1,\ldots,u_m]]}
 {(H_1(c+u),\ldots,H_m(c+u))},
 \qquad c_\alpha=p_{t_\alpha}.
\]
This explains why simply counting compatible partitions
does not compute their local multiplicities.

\section{Which labels and partitions can contribute?}

A necessary condition for $m_{A,B;\pi}(p)\ne0$ is
\begin{equation}\label{eq:compatible}
 A\subseteq A_0,\qquad B\subseteq B_0,\qquad
 \pi\leq \sigma_{A,B}(p),
\end{equation}
where $\sigma_{A,B}(p)$ is the equality partition of the remaining
coordinates and $\leq$ denotes refinement.
Explicitly, its blocks are $A_0\setminus A$, $B_0\setminus B$
when nonempty, and the singleton ordinary labels.
No block of $\pi$ can mix these groups.

For balanced choices $|A|=|B|=s\leq k$, compatibility is also
sufficient for $p$ to solve the auxiliary equations.
If $s<k$, the remaining opposite singular cluster makes each
singular reduced equation vanish; at an ordinary coordinate,
\eqref{eq:reduced} specializes to
$(z_u^2-1)^{k-s}F_u(z)=0$.
For $s=k$, only \eqref{eq:fully-reduced} remains.
These statements give membership, not the value of the local length.

Under Definition~\ref{def:simple}, fixing all $k$ pairs leaves just
one compatible choice:
\[
 A=A_0,\qquad B=B_0,\qquad \pi=\hat0_{R_0},
 \qquad m_{A_0,B_0;\hat0_{R_0}}(p)=1.
\]
For $s<k$, the multiplicities remain the local algebra lengths
in \eqref{eq:local-length}; we impose no formula for them.

\section{Depermutation weights}

For a partition $\pi$ define its partition-lattice weight by
\[
 \mu(\hat0,\pi)=\prod_{S\in\pi}(-1)^{|S|-1}(|S|-1)!.
\]
For balanced $A,B$, put $s=|A|=|B|$ and $r_s=N-2s$. Define
\[
 Z^{\mathrm{dep}}_{A,B}
 =\frac1{r_s!}\sum_{\pi\in\Pi(R)}
                  \mu(\hat0,\pi)Z_{A,B;\pi},
\]
and then
\[
 Z_s^{\mathrm{dep}}
 =\frac{(s!)^2r_s!}{N!}
   \sum_{\substack{A\cap B=\varnothing\\|A|=|B|=s}}
     Z^{\mathrm{dep}}_{A,B}.
\]
Here $\Pi(R)$ is the set of partitions of $R$.
For an empty set there is one partition, with weight $1$.
These are rational weighted cycles; their coefficients need not be
individual local lengths.

Let $w_s(p)=\coeff_p Z_s^{\mathrm{dep}}$. Cancelling the two
$r_s!$ factors gives the explicit definition
\begin{equation}\label{eq:weight}
 \boxed{
 w_s(p)=\frac{(s!)^2}{N!}
 \sum_{\substack{A\cap B=\varnothing\\|A|=|B|=s}}
 \ \sum_{\pi\in\Pi([N]\setminus(A\cup B))}
       \mu(\hat0,\pi)m_{A,B;\pi}(p).}
\end{equation}
Only the compatible terms in \eqref{eq:compatible} contribute.
In particular, $w_s(p)=0$ for $s>k$, and simplicity gives
\begin{equation}\label{eq:top}
 w_k(p)=\frac{(k!)^2}{N!}.
\end{equation}

For bookkeeping, choose one compatible pair $A_s\subseteq A_0$,
$B_s\subseteq B_0$ of size $s$ and set
\[
 M_s(p)=\sum_{\pi\leq\sigma_{A_s,B_s}(p)}
          \mu(\hat0,\pi)m_{A_s,B_s;\pi}(p).
\]
Permuting the $k$ plus labels and the $k$ minus labels fixes $p$,
so this sum is independent of the choice of $A_s,B_s$.
There are $\binom ks^2$ such choices. Consequently
\begin{equation}\label{eq:Ms}
 w_s(p)=\frac{(s!)^2\binom ks^2}{N!}M_s(p)
       =\frac{(k!)^2}{N!((k-s)!)^2}M_s(p),
 \qquad M_k(p)=1.
\end{equation}
This is a symmetry and normalization identity, not a determination
of the unknown lengths inside $M_s(p)$.

\section{The alternating local index}

Define the candidate corrected cycle by
\[
 \widetilde Z=
 \sum_{s=0}^{\lfloor N/2\rfloor}
        (-1)^s\binom Ls Z_s^{\mathrm{dep}}.
\]
\begin{definition}[Local index]
The index at the ordered configuration $p$ is its coefficient in this
cycle:
\begin{equation}\label{eq:index}
 \boxed{
 \ind_p=\coeff_p\widetilde Z
       =\sum_{s=0}^k(-1)^s\binom Ls w_s(p)
       =\frac{(k!)^2}{N!}
          \sum_{s=0}^k
          \frac{(-1)^s\binom Ls}{((k-s)!)^2}M_s(p).}
\end{equation}
\end{definition}

For $k=1$, every compatible partition is the singleton partition.
Writing $\ell_{\mathrm{full}}(p)$ for the local length of
$(P_1,\ldots,P_N)$, the definition reduces to
\[
 M_0(p)=\ell_{\mathrm{full}}(p),\qquad M_1(p)=1,\qquad
 \ind_p=\frac{\ell_{\mathrm{full}}(p)-L}{N!}.
\]
Even here, reduced simplicity alone is not a definition of
$\ell_{\mathrm{full}}(p)$.

For $k=2$, take $A_0=\{1,2\}$, $B_0=\{3,4\}$.
With no fixed labels, let $\ell_{00}$ denote the full local length,
$\ell_{10}$ the length with $\{1,2\}$ merged,
$\ell_{01}$ the length with $\{3,4\}$ merged, and
$\ell_{11}$ the length with both blocks merged.
All other blocks are singletons. With one fixed pair, let
$\ell_{\mathrm{one}}$ be the local length for any one compatible
choice and the singleton partition. Then
\[
 M_0=\ell_{00}-\ell_{10}-\ell_{01}+\ell_{11},\qquad
 M_1=\ell_{\mathrm{one}},\qquad M_2=1,
\]
so
\[
 \ind_p=
 \frac{\ell_{00}-\ell_{10}-\ell_{01}+\ell_{11}
       -4L\ell_{\mathrm{one}}+4\binom L2}{N!}.
\]
The factors $4$ come from the fixed-label normalization and the
four choices of one plus label and one minus label.

The local question is therefore to determine the lengths
$m_{A,B;\pi}(p)$, or the combinations $M_s(p)$, under the fully
reduced simplicity assumption. Formula~\eqref{eq:index} defines
the index but does not establish its positivity, a component-degree
interpretation, or its expected vanishing for $k\geq2$.
\end{document}
